A ternary node has three sites that could carry a number format; one needs to.
The weight is a code: multiplying by it is a sign-select. The sample arrives
ADC-native. The accumulator alone has a range to spend, and Ternary Network
Floats spend it --- a sign, a balanced-ternary exponent of $E_t$ trits, an
$M$-bit mantissa, $1+E_t+M=N$.

\paragraph{Budget convention, stated once and binding throughout.}\label{sec:budget}
$1+E_t+M=N$ counts the exponent as $E_t$ \emph{trit positions}. It is a
specification convention, not a storage claim, and this paper uses two budgets
under two explicit names. The \textbf{nominal} budget is $1+E_t+M=N$; the
\textbf{physical} budget is $1+\lceil E_t\log_2 3\rceil+M \le N$, which is what a
packed binary word costs. The two differ by
$\lceil E_t\log_2 3\rceil - E_t$ cells, three at $E_t=4$ and seven at $E_t=11$, so
the adopted 16-bit rung $(E_t,M)=(4,11)$ occupies $19$ physical bits.
\textbf{The canonical name of a rung is now its bit width}: TNF4, TNF8, TNF16,
TNF32 and TNF64 are 6, 10, 19, 36 and 69 bits, and the nominal names are retained
only as aliases, because the sequence $4,8,16,32,64$ is exactly the IEEE width
sequence and a row reading ``TNF16 against binary16'' looks matched and is not.
Every accuracy comparison against a binary format is stated with its budget named.
Under the physical budget TNF holds no advantage at $16$ or $32$ bits ---
Section~\ref{sec:threeresults} measures this, Table~\ref{tab:accphys} states it
for the accuracy table specifically, and Theorem~\ref{thm:rangeperbit} says why it
cannot be otherwise. Claims of the form ``TNF wins at equal width'' appear nowhere
in this paper without the budget attached.

The weight is \textbf{GFTernary}, a two-bit $\{-\varphi,0,+\varphi\}$ alphabet,
and it is where the multiplier is decided. Removing one rather than cheapening
it requires the weight lattice to be \emph{closed} under weight application and
accumulation alike. $\mathbb{Z}[\varphi]$ is, because $\varphi^2=\varphi+1$:
applying a weight is the Fibonacci step $(a,b)\mapsto(b,a+b)$, accumulation is
componentwise, and a layer's linear path is exact. Machine-checked in Coq.

Closure does not single out $\varphi$; enumeration does. A scale is
multiplier-free exactly when it is an algebraic integer whose companion matrix
has entries in $\{0,\pm1\}$. Degree two admits $\varphi$ alone, and nothing lies
between the shift and $\varphi$. The family $r^d=r+1$ reaches any granularity at
\emph{one addition} and $d$ registers: fineness costs registers, not adders.

\textbf{The throughput-per-area ranking is withdrawn.} It read
$0.1797$\,MHz/LUT for GFTernary against $0.1631$ for the next entry, and it
cannot stand: re-measurement found that on-die check logic had been folded to
constants at synthesis in fourteen of twenty-eight clauses, so every area figure
in that comparison measured a design with its checks deleted, and no honest
(area, frequency) pair existed anywhere in the corpus. Taken again with both
numbers from the same build, frequency follows
$F_{\max}=4397\cdot\mathrm{LUT}^{-0.648}$ ($R^2=0.920$) across a
$39.6\times$ span of area, which makes MHz per LUT a proxy for size rather than a
figure of merit, and the residuals span a factor of two, so size explains $92\%$
of frequency and structure the rest. What survives is an \emph{area} result,
which no re-measurement moved: on an XC7A200T the folded ternary neuron costs
$28.0$ LUT per weight at fan-in $8$, rising to $33.1$ and $33.2$ at fan-in $16$
and $32$, with \emph{no DSP at any fan-in} (Table~\ref{tab:node}). Formats
needing an external learned scale are tabulated apart.

Which rung follows from the budget --- $\varphi$ at four bits, $r^5{=}r^3{+}1$ at
five, $r^6{=}r{+}1$ at six --- reaching $3.7\%$ of \texttt{fp32} at six bits and
one adder. On the block axis the element is closed, the best eight-level
codebook beating a deployed one by $0.9\%$, and the scale is not: a $\varphi$
grid at four bits yielded lower perplexity than MXFP4 in the two evaluated models,
at $0.125$ fewer bits per weight.

On a physically ternary store the block comparison changes sign. A 27-level
element in one three-trit cell, with a power-of-two scale per $32$ weights,
gives lower perplexity than NVFP4, MX+ and MXFP4 on SmolLM2-135M ($16.29$
against $17.47$, $17.98$ and $21.94$) in $101$ trits per $32$ weights against
$106$, $106$ and $102$, in a pre-registered run (Section~\ref{sec:t27block}).
Counted in bits it costs $161$ against NVFP4's $144$ and sits on the line
between two binary formats: it is the trit slack of
Corollary~\ref{cor:slackartefact}, measured, and it pays only on a native
ternary cell costing under $1.43$ bit cells.

Applicability is stated as a threshold and then tested against it. The scalar is
$D=\operatorname{mean}(\lvert\log_2|x|\rvert)$; the threshold is a property of the
rung, not of the family --- $D \gtrsim 9.5$ at sixteen physical cells and
$D \gtrsim 21.9$ at thirty-two, with a third rung admitting no threshold at all ---
and it is a conjunction with the range condition, which does most of the excluding.
Eleven intermediate quantities of a real GPT-2 block satisfy none of it, because
normalisation is what destroys the property. Seven of twenty-two candidate pairs
outside neural inference do satisfy it, at $1.23\times$ to $4.00\times$: raw-SI
inverse-square denominators, unnormalised likelihood products, importance weights.
\texttt{takum} reaches $\pm 255$ binades at both widths compared and never leaves
range where TNF does, so the format that wins inside the window is the one that has
to be selected into it. One task outcome rather than a round-trip error is also
reported: on a maximum-a-posteriori estimate in raw SI units with a fixed
\texttt{float32} accumulator, TNF32 lands $1.60\times$ closer to the
\texttt{float64} solution than \texttt{takum32} and settles in $21$ iterations
against $27$, while \texttt{binary32} beats both by a factor of thirty-three.

\textbf{At matched physical width the comparison is negative.} Computed from
the committed reference implementations, posit holds more representable values
than TNF and a finer step at unity at every width the ladder occupies --- $12\times$
finer at 19 bits, $4\times$ at 10, $2\times$ at 6 --- and posit with $es{=}2$ weakly
dominates on both coordinates at all three. A ternary exponent does not tile a
binary field: four trits use $81$ of $128$ codes, leaving $8{,}190$ of the $2^{19}$
words unreachable. The no-survivors corollary is accordingly withdrawn to a
matched-\emph{name} claim (Section~\ref{sec:matchedwidth}). The case for this
format was never precision per bit; it is multiplier removal, decode cost and the
rungs where nothing else trains.

Section~\ref{sec:selfcheck} sets out the ways a format comparison fails while
each of its measurements is correct, since every quantitative claim above depends
on those failures having been ruled out.
